\section*{Bab \ref{ch:b3:renal-osmoregulation} --- Fisiologi Ginjal dan Osmoregulasi}

\begin{solution}{exo:b3:renal-osmoregulation:1}
\omterm{def:b3:renal-osmoregulation:nephron}{Glomerulus}: menapis plasma dikurangi proteinnya. \omterm{def:b3:renal-osmoregulation:nephron}{Tubulus proksimal}: menyerap kembali
dua pertiga garam dan airnya, semua glukosa dan asam aminonya,
serta bikarbonatnya. Cabang turun Henle: kehilangan air ke medulanya.
Cabang naik: memompa keluar NaCl, tidak melewatkan air. Tubulus distal:
menyesuaikan natrium (\omterm{def:b3:renal-osmoregulation:controls}{aldosteron}), menyekresikan kalium. \omterm{def:b3:renal-osmoregulation:nephron}{Saluran pengumpul}:
penyerapan kembali air di bawah vasopresin, penanganan proton dan urea.
\end{solution}

\begin{solution}{exo:b3:renal-osmoregulation:2}
Bersihan: volume plasma yang dibersihkan dari zatnya tiap satuan
waktu, $UV/P$. Di sini $100\times 1.5/2 = \qty{75}{mL/min}$, di bawah LFG
\qty{120}{mL/min}: zatnya ditapis lalu sebagian diserap kembali
(sekitar \qty{38}{\%} jumlah yang ditapis).
\end{solution}

\begin{solution}{exo:b3:renal-osmoregulation:3}
Seekor ikan air tawar lebih asin daripada mediumnya: airnya masuk lewat
insangnya secara osmosis, sehingga minum hanya akan menambah banjirnya; ia
mengekskresikan kemih encer yang berlimpah dan \omterm{def:b3:renal-osmoregulation:comparative}{sel klorida} insangnya mengambil
Na$^{+}$ dan Cl$^{-}$ secara aktif dari airnya. Seekor ikan laut kurang
asin daripada lautnya: ia kehilangan air lewat insangnya, meminum air laut
untuk menggantinya, dan \omterm{def:b3:renal-osmoregulation:comparative}{sel kloridanya} memompa garam yang ditelannya keluar.
\end{solution}

\begin{solution}{exo:b3:renal-osmoregulation:4}
\omterm{prop:b3:renal-osmoregulation:countercurrent}{Efek tunggalnya} adalah selisih \qty{200}{mOsm/L} yang dipertahankan cabang
naik yang tebal antara interstisiumnya dan cairannya sendiri
dengan memompa NaCl keluar sambil tidak membiarkan air mengikutinya. Cabang naiknya
harus tidak melewatkan air, atau airnya akan mengikuti garamnya dan
selisihnya akan lenyap; cabang turunnya harus melewatkan air sehingga
cairan yang mencapai tikungannya sudah dipekatkan oleh
interstisiumnya dan tiap tahapnya bekerja pada cairan yang lebih asin daripada yang sebelumnya. Dua
cabang dengan keterlewatan yang sama tidak dapat melipatgandakan.
\end{solution}

\begin{solution}{exo:b3:renal-osmoregulation:5}
Ujung aferennya: $60 - 18 - 28 = \qty{14}{mmHg}$; ujung eferennya: $60 - 18 -
35 = \qty{7}{mmHg}$. Ketika airnya ditapis, protein yang tertinggal
memekat dan tekanan onkotiknya naik sepanjang kapilernya; bila ia
mencapai \qty{42}{mmHg}, tekanan netonya nol dan penapisannya berhenti
sebelum ujungnya (keseimbangan penapisan). LFG-nya lalu bergantung pada
aliran plasmanya: aliran yang lebih cepat memekatkan proteinnya lebih lambat, sehingga
penapisannya berlanjut sepanjang bentangan yang lebih panjang dan LFG-nya naik seiring
alirannya.
\end{solution}

\begin{solution}{exo:b3:renal-osmoregulation:6}
Muatan yang ditapis: $\qty{40}{mL/min}\times\qty{0.14}{mmol/mL} =
\qty{5.6}{mmol/min}$. Ekskresinya: $70\times 1 = \qty{0.07}{mmol/min}$.
Yang diserap kembali: $5.53/5.6 = \qty{98.75}{\%}$. Bersihan natriumnya: $70\times
1/140 = \qty{0.5}{mL/min}$.
\end{solution}

\begin{solution}{exo:b3:renal-osmoregulation:7}
$900/1200 = \qty{0.75}{L}$ tiap hari. Pada \qty{400}{mOsm/L}: $900/400 =
\qty{2.25}{L}$ kemih, sehingga sekitar $2.25 + 0.9 = \qty{3.15}{L}$ harus
diminum (dikurangi air di dalam makanannya): pasien dengan mekanisme
pemekatan yang gagal harus minum, dan ia terancam bila tidak dapat.
\end{solution}

\begin{solution}{exo:b3:renal-osmoregulation:8}
(a) Diabetes insipidus: \qtyrange{10}{15}{L} kemih pada
\qtyrange{50}{100}{mOsm/L}, natrium plasmanya tinggi kapan pun minumnya tertinggal,
haus tak henti. (b) Vasopresin yang tidak pada tempatnya: kemih pekat yang sedikit,
airnya ditahan, natrium plasmanya rendah karena pengenceran, hausnya tidak
bertambah meskipun minumnya berlanjut --- kebingungan dan kejang ketika
natriumnya turun jauh. (c) Kelebihan \omterm{def:b3:renal-osmoregulation:controls}{aldosteron}: natriumnya ditahan bersama airnya,
volumenya mengembang, tekanan darahnya tinggi, kaliumnya hilang (K$^{+}$ plasmanya
rendah), natrium plasmanya nyaris normal sebab airnya mengikuti garamnya dan
ginjalnya lolos sebagian; volume kemihnya normal. (d) Makanan tanpa garam: renin
dan \omterm{def:b3:renal-osmoregulation:controls}{aldosteronnya} naik dalam sehari, natrium kemihnya turun menjadi beberapa
milimol tiap hari, natrium plasmanya tertahan, volume dan tekanannya sedikit
lebih rendah.
\end{solution}

\begin{solution}{exo:b3:renal-osmoregulation:9}
Kemih encer: bersihan osmolarnya $100\times 6/300 = \qty{2}{mL/min}$;
\omterm{prop:b3:renal-osmoregulation:water}{bersihan air bebasnya} $6 - 2 = +\qty{4}{mL/min}$: ginjalnya sedang membuang
air bebas zat terlarut, seperti sesudah sebuah muatan air dengan vasopresinnya dimatikan.
Kemih pekat: bersihan osmolarnya $1000\times 0.6/300 =
\qty{2}{mL/min}$; \omterm{prop:b3:renal-osmoregulation:water}{bersihan air bebasnya} $0.6 - 2 = -\qty{1.4}{mL/min}$:
\qty{1.4}{mL/min} air bebas sedang dikembalikan ke tubuhnya ---
yakni dehidrasi, dengan vasopresinnya tinggi.
\end{solution}

\begin{solution}{exo:b3:renal-osmoregulation:10}
Nomorilah tahapnya dari korteksnya; misalkanlah interstisium pada tahap $k$
adalah $I_{k}$. Cairan yang turun pada $k$ sama dengan $I_{k}$; cairan yang
naik pada $k$ adalah $I_{k} - 200$; dan cairan yang naik pada $k$ adalah apa yang
meninggalkan tikungannya lalu melewati tahap $n, n-1, \dots, k+1$. Pada keadaan mantap,
cairan yang masuk pada \qty{300}{} meninggalkan puncak cabang naiknya
pada $I_{1} - 200$, dan interstisium tiap tahapnya melebihi yang di atasnya
sampai sebesar \omterm{prop:b3:renal-osmoregulation:countercurrent}{efek tunggalnya}, sehingga $I_{k} \le 300 + 200k$ dengan kesamaan ketika
tiap tahapnya bekerja pada efek penuh: tikungannya mencapai $300 + 200n$. Landaiannya
karena itu dibatasi panjang lengkungnya (yakni cacah
tahapnya) kali \omterm{prop:b3:renal-osmoregulation:countercurrent}{efek tunggalnya}, bukan oleh daya pompa sel mana pun. $n$
adalah panjang lengkungnya dibagi jarak yang sepanjangnya cairannya
menyeimbang dengan interstisiumnya: lengkung jukstamedula yang panjang, dan
lengkung yang sangat panjang milik hewan pengerat gurun, memberi $n$ yang besar; kapasitas pompanya
dan penghanyutan zat terlarut oleh vasa rektanya membatasi efeknya.
\end{solution}

\begin{solution}{exo:b3:renal-osmoregulation:11}
\omterm{def:b3:renal-osmoregulation:comparative}{Air metabolisme}: $5\times 0.8 = \qty{4}{g}$ pati memberi $4\times 0.6
= \qty{2.4}{g}$. Kehilangannya: kemih $3/5500\ \unit{L} = \qty{0.55}{mL}$;
penguapan \qty{1.5}{g} (yang sudah neto dari pemulihan hidungnya). Total
kehilangannya \qty{2.05}{g} terhadap perolehan \qty{2.4}{g}: sebuah lebihan
\qty{0.35}{g}, yang menutupi kehilangan tinjanya yang kecil dan air higroskopis
bijinya menjadi bonus. Ia bertahan hidup tanpa minum. Tanpa
pemulihan hidungnya, kehilangan napasnya akan menjadi $1.5/0.3 = \qty{5}{g}$,
dan ia akan mati dalam beberapa hari.
\end{solution}

\begin{solution}{exo:b3:renal-osmoregulation:12}
Pada aliran sejajar, darah dan dialisatnya menyeimbang di tengah jalan sepanjang
selaputnya lalu selisih kepekatan yang mendorong pembauran turun
menjadi nol: paling baik pun darahnya pergi dengan separuh ureanya. Pada aliran arus
balik, dialisat tersegar menjumpai darah terbersih dan
landaiannya dipertahankan sepanjang seluruh selaputnya, sehingga darahnya dapat
pergi dengan nyaris kepekatan masuk dialisatnya --- yakni asas
yang sama dengan \omterm{def:b3:renal-osmoregulation:nephron}{lengkung Henle}. Bersihannya: $300\times 0.7 =
\qty{210}{mL/min}$ selama $3\times 4 = \qty{12}{h}$ seminggu, yakni
$210\times 720 = \qty{151}{L}$ seminggu; dua ginjal pada
\qty{125}{mL/min} membersihkan $125\times\num{10080} = \qty{1260}{L}$.
Mesinnya mengerjakan sekitar seperdelapan kerja ginjalnya, yakni rata-rata waktu
\qty{15}{mL/min} --- cukup untuk hidup, tidak cukup untuk sehat.
\end{solution}

\begin{solution}{pb:b3:renal-osmoregulation:1}
\textbf{1.} Tekanan netonya $55 - 15 - 30 = \qty{10}{mmHg}$; LFG $= 12.5
\times 10 = \qty{125}{mL/min}$.
\textbf{2.} $125\times 1440 = \qty{180}{L/d}$; $180/3 = 60$ kali.
\textbf{3.} Inulin: $62.5\times 1/0.5 = \qty{125}{mL/min}$, sama dengan
LFG-nya. Kreatinin: $1.3\times 1/0.01 = \qty{130}{mL/min}$, sedikit
lebih besar sebab \omterm{def:b3:renal-osmoregulation:nephron}{tubulus proksimalnya} menyekresikan sejumlah kecil
kreatinin di samping yang ditapis.
\textbf{4.} Bersihan PAH $13\times 1/0.02 = \qty{650}{mL/min}$, yakni
aliran plasma ginjalnya; aliran darahnya $650/(1 - 0.45) = \qty{1180}{mL/min}$,
sekitar \qty{1.2}{L/min}; pecahan penapisannya $125/650 = 0.19$.
\textbf{5.} Tekanan netonya \qty{15}{mmHg}, LFG-nya \qty{188}{mL/min} (lebih kecil pada
kenyataannya, sebab tekanan onkotiknya naik sepanjang kapilernya).
Angiotensin II menyempitkan arteriol eferennya lalu menaikkan
tekanan \omterm{def:b3:renal-osmoregulation:nephron}{glomerulusnya}; menghalanginya menurunkan tekanan dan LFG-nya
sedikit. Pada diabetes, \omterm{def:b3:renal-osmoregulation:nephron}{glomerulusnya} menapis di bawah tekanan tinggi dan
tapisannya aus; menurunkan tekanannya memperlambat kerusakannya bertahun-tahun.
\textbf{6.} Podosit dan diafragma celahnya (atau muatan
selaput dasarnya): yakni lapisan yang menahan albuminnya. Kehilangan
\qty{3}{g/d} menurunkan tekanan onkotik plasmanya; cairannya meninggalkan
kapilernya di mana-mana dan jaringannya membengkak --- yakni sembab
sindrom nefrotiknya.
\textbf{7.} Yang ditapis $125\times 5 = \qty{0.625}{mmol/min}$; yang diekskresikan
$250\times 1 = \qty{0.25}{mmol/min}$; yang diserap kembali $0.375$, yakni
\qty{60}{\%}; bersihannya $250/5 = \qty{50}{mL/min}$.
\textbf{8.} Yang ditapis $180\times 0.14 = \qty{25.2}{mol/d}$ Na$^{+}$,
yakni \qty{580}{g} natrium, \qty{1.47}{kg} sebagai NaCl. Yang diekskresikan $100\times 1.44
= \qty{144}{mmol/d}$; yang diserap kembali $1 - 144/\num{25200} = \qty{99.4}{\%}$.
Nyaris satu setengah kilogram garam akan hilang tiap hari.
\textbf{9.} Yang ditapis $125\times 5 = \qty{0.625}{mmol/min}$, yakni
\qty{112}{mg/min}, \qty{162}{g/d}; tidak ada yang diekskresikan, sebab ini di bawah
$T_{m}$. Muatan yang ditapis mencapai $T_{m}$ pada $2.1/0.125 =
\qty{16.8}{mmol/L}$ (\qty{300}{mg/dL}) --- lebih tinggi daripada ambang
yang teramati sekitar \qty{10}{mmol/L}, sebab \omterm{def:b3:renal-osmoregulation:nephron}{nefronnya} berbeda-beda dan
yang terlemah tumpah lebih dulu (yakni pemencaran kurva titrasinya).
\textbf{10.} Yang ditapis $125\times 20 = \qty{2.5}{mmol/min}$; yang diekskresikan
$2.5 - 2.1 = \qty{0.4}{mmol/min}$, \qty{72}{mg/min}, yakni $576$ mmol atau
\qty{104}{g} tiap hari. Kemih tambahannya $576/300 = \qty{1.9}{L}$. Glukosanya
menyeret air ke dalam kemihnya (diuresis osmotik), pasiennya mengalami dehidrasi
dan haus; \qty{104}{g} glukosa adalah \qty{1.8}{MJ} yang hilang tiap hari, dan
selnya, yang tidak dapat mengambil glukosa, membakar lemak dan protein: bobotnya
turun.
\textbf{11.} Yang diserap kembali $\approx \qty{25}{mol/d}$; ATP-nya $25/3 =
\qty{8.4}{mol/d}$; $8.4\times 50 = \qty{420}{kJ}$, yakni sekitar \qty{6}{\%}
dari \qty{7}{MJ}.
\textbf{12.} Untuk menyekresikan tiap limbahnya, tubulusnya akan memerlukan sebuah pengangkut
yang mengenalinya --- ribuan di antaranya, dan tidak satu pun bagi sebuah molekul yang tak pernah
dijumpai sebelumnya. Menapis semua yang kecil lalu menyerap kembali beberapa ratus
jenis yang diinginkan tubuhnya hanya memerlukan pengangkut bagi apa yang diketahui
berharga; apa pun yang tak dikenali pergi secara bawaan. Harganya adalah
tenaga pertanyaan~11.
\textbf{13.} Bersihan osmolarnya $600\times 1/290 = \qty{2.07}{mL/min}$;
\omterm{prop:b3:renal-osmoregulation:water}{bersihan air bebasnya} $1 - 2.07 = -\qty{1.07}{mL/min}$: ginjalnya sedang
menghemat air.
\textbf{14.} Minimumnya $600/1200 = \qty{0.5}{L}$; maksimumnya $600/50 =
\qty{12}{L}$; kebutuhan hariannya pada minimumnya $0.5 + 0.9 = \qty{1.4}{L}$.
\textbf{15.} \qty{1000}{mOsm} pada \qty{1200}{mOsm/L} memerlukan
\qty{0.83}{L} kemih: perolehan netonya \qty{0.17}{L} sebelum zat terlarut hari itu.
Sesudahnya: zat terlarut total hari itu \qty{1600}{mOsm}, yakni \qty{1.33}{L}
kemih, ditambah \qty{0.9}{L} tak terasa, yakni \qty{2.23}{L} keluar bagi \qty{1}{L}
masuk --- yakni sebuah kekurangan \qty{1.23}{L}, terhadap \qty{1.4}{L} tanpa minum.
Seliter air laut itu memperbaiki neracanya hanya sebesar \qty{0.17}{L}.
\textbf{16.} Kemih pada \qty{50}{mOsm/L} bagi hari itu: $600/50 =
\qty{12}{L}$; \qty{12.9}{L} harus diminum tiap hari, yakni sekitar \qty{13}{L}.
\textbf{17.} $B$ liter bir membawa $20B$ mOsm; kemih maksimalnya adalah
$(600 + 20B)/50 = 12 + 0.4B$ liter, dan ia harus mengusung $B - 0.9$
liter: $B - 0.9 \le 12 + 0.4B$, $B \le \qty{21.5}{L}$. Sekitar dua puluh
liter --- tetapi bila sedikit saja yang lain dimakan, zat terlarut hari itu turun menjadi
\qty{200}{mOsm} dan batasnya jatuh ke sekitar \qty{8}{L}: yakni
hiponatremia peminum berat yang tidak makan.
\textbf{18.} $298/290 - 1 = \qty{2.8}{\%}$: jauh melewati \qty{1}{\%} yang
memicu vasopresin dan haus; keduanya maksimal. Dengan zat terlarut tetap:
$290\times 42 = 298\times V$, $V = \qty{40.9}{L}$, yakni sebuah kekurangan sekitar
\qty{1.1}{L}.
\textbf{19.} Enam liter air mengencerkan plasmanya sementara keringatnya menyingkirkan
garamnya; lalu olahraga, nyeri, dan tekanan melepaskan vasopresin tanpa memedulikan
\omterm{def:b3:renal-osmoregulation:problem}{osmolaritasnya}, sehingga ginjalnya tidak dapat mengencerkan kemihnya cukup untuk mengekskresikan
lebihannya. Natrium plasmanya turun, sel otaknya membengkak --- kejang dan kematian
telah menyusul. Ginjalnya akan memerlukan vasopresin yang dimatikan dan
\qty{6}{L} kemih pada $U_{\min}$; pelarinya sebaiknya minum menurut hausnya lalu
mengganti garamnya.
\textbf{20.} $3/5500\ \unit{L} = \qty{0.55}{mL}$ tiap hari. Sebuah ginjal manusia
akan memerlukan $3/1200 = \qty{2.5}{mL}$ bagi zat terlarut yang sama: tikus
kangurunya memakai seperlima airnya.
\textbf{21.} Masuk: \qty{2.4}{g}. Keluar: $1.6 + 0.3 + 0.55 = \qty{2.45}{g}$.
Seimbang sampai dalam \qty{0.05}{g} tiap hari, yang dipasok air higroskopis
bijinya (beberapa persen dari \qty{5}{g}).
\textbf{22.} Airnya $500\times 0.65 = \qty{325}{L}$; seperempatnya
\qty{81}{L}; pada \qty{5}{L/d}, yakni enam belas hari.
\textbf{23.} $150/800 = \qty{0.19}{L}$, yakni sekitar \qty{190}{mL} air garam
--- kurang daripada \qty{100}{mL} air laut dan air ikannya
bersama-sama, sehingga burungnya memperoleh air. Ginjalnya, pada \qty{600}{mOsm/L}
(\qty{300}{mmol/L} NaCl), akan memerlukan \qty{0.5}{L} kemih untuk mengusung
garam yang sama, yakni lebih banyak air daripada yang diambilnya: sebuah ginjal yang tidak dapat mengalahkan
air laut tidak dapat membuat seekor burung meminumnya, sedangkan kelenjarnya, pada dua kali
kepekatan lautnya, dapat.
\textbf{24.} \qty{30}{g} air masuk, sehingga sekitar \qty{30}{mL} kemih tiap
hari, pada sepersekian kecil \omterm{def:b3:renal-osmoregulation:problem}{osmolaritas} darahnya (beberapa puluh
mOsm/L terhadap \num{300}). Garamnya pergi di dalam kemih itu dan lewat pembauran
melintasi insangnya; tanpa pengambilan aktif oleh \omterm{def:b3:renal-osmoregulation:comparative}{sel kloridanya} dari
air yang memuat semilimol tiap liter, ikannya akan terkuras dalam
beberapa hari.
\textbf{25.} LFG \qty{125}{mL/min}, yakni \qty{180}{L} tiap hari; kemih
wajibnya \qty{0.5}{L} tiap hari; seliter air laut menghasilkan neto
\qty{0.17}{L} begitu garamnya diekskresikan --- yakni nyaris tidak ada.
\end{solution}
